\section{Билет 69. Микро и макросостояние термодинамической системы. Статистический вес и вероятность макросостояния. Статистическое определение энтропии. Формула Больцмана.}

\begin{definition}
Термодинамическую систему, состоящую из огромного числа частиц ($N \sim 10^{23}$), можно описывать на двух уровнях:
\begin{itemize}
    \item \textbf{Макросостояние} характеризуется набором макроскопических параметров, доступных для непосредственного измерения: давлением $p$, объёмом $V$, температурой $T$, внутренней энергией $U$ и т.д.
    \item \textbf{Микросостояние} задаётся точным указанием координат $\vec{r}_i$ и импульсов $\vec{p}_i$ (или скоростей $\vec{v}_i$) каждой из $N$ частиц системы. Одно и то же макросостояние может быть реализовано колоссальным числом различных микросостояний.
\end{itemize}
\end{definition}

\begin{definition}
\textbf{Статистическим весом} (термодинамической вероятностью) $W$ макросостояния называется число различных микросостояний, которые реализуют данное макросостояние.
\par\noindent\textbf{Вероятность} реализации макросостояния пропорциональна его статистическому весу:
\begin{equation}
    P \propto W.
\end{equation}
Система стремится перейти в состояние с наибольшим $W$, так как оно является наиболее вероятным. Равновесное состояние соответствует максимуму статистического веса: $W = W_{\max}$.
\end{definition}

\begin{theorem}[Статистическое определение энтропии]
Энтропия $S$ является мерой числа способов, которыми может быть реализовано данное макросостояние системы. Статистическое определение энтропии даётся \textbf{формулой Больцмана}:
\begin{equation}
    S = k \ln W, \label{eq:boltzmann_entropy_69}
\end{equation}
где $k = 1.38 \cdot 10^{-23}$ Дж/К --- постоянная Больцмана, $W$ --- статистический вес макросостояния.
\end{theorem}

\begin{property}[Почему логарифм?]
Функция энтропии должна быть \textbf{аддитивной}: при объединении двух независимых подсистем 1 и 2 энтропия всей системы равна сумме энтропий: $S_{12} = S_1 + S_2$.
Статистический вес, напротив, \textbf{мультипликативен}: число комбинаций состояний независимых подсистем перемножается: $W_{12} = W_1 \cdot W_2$.
Единственная функция, превращающая произведение в сумму --- логарифм. Поэтому $S \sim \ln W$, а коэффициент пропорциональности выбирается равным $k$ для согласования с термодинамической шкалой.
\end{property}

\begin{remark}
\textbf{Связь со вторым началом термодинамики.}
\begin{itemize}
    \item В изолированной системе самопроизвольные процессы идут в направлении увеличения статистического веса $W$, а значит, и энтропии $S$: $\Delta S \ge 0$.
    \item Равновесное состояние --- это состояние с максимальным $W$ и максимальной $S$.
    \item Формула Больцмана раскрывает статистическую сущность второго начала: необратимость процессов обусловлена колоссальным различием в числе микросостояний для упорядоченных и хаотических конфигураций. Переход от порядка к хаосу статистически неизбежен, хотя в принципе обратный переход не запрещён законами механики (его вероятность ничтожно мала).
\end{itemize}
\end{remark}
